\section{Conclusion and Future Work}

We presented Polygraph, a tool that classifies \ac{sast} findings as \acp{tp}
or \acp{fp} from a compact, finding-specific context. It reads and writes
standard \ac{sarif} and records its verdicts in the standard
\texttt{suppressions} field, so it operates between any compliant scanner and any
compliant viewer. Its two-stage design keeps the common case cheap, as only the
findings the static context does not settle escalate to an autonomous agent that
explores the repository on its own.

On the OWASP Benchmark for Python, the best of four open-weight models reached an
F1-score of 0.977 on the \ac{fp} class, suppressing 95.8\% of the traps at
a dangerous-suppression rate of 0.4\%. On RealVuln, whose verdicts depend on code
spread across a repository, it reached an F1-score of 0.797 and suppressed 82.5\%
of the traps at a dangerous-suppression rate of 5.0\%. Escalation remained rare
where the static context sufficed, firing on 3.6\% of the OWASP findings, and
recovered 145 real vulnerabilities on RealVuln for 7 retained traps. Parameter
count did not predict quality, as the 27B model led and the 80B model ranked
last.

Several directions remain open. All runs ran without data-flow analysis, since it
is the slowest step of the pipeline, so the contribution of this context signal
is still unmeasured. Selecting the enrichers per finding,
guided by its \ac{cwe} class, would limit that cost to the findings whose
verdict depends on it. Second,
the evaluation covers Python alone and consumes synthesized reports rather than
live scanner output, so a broader study across languages, scanners, and
production code is needed to show how far these results transfer. Finally, given that
model size did not decide quality, post-training a small open-weight model
with confirmed real-world \acp{tp} and \acp{fp} could result in better results overall.
